% v0_3_theorems.tex
% Framework Theorems (working draft)
%
% This file is intentionally NOT included from main.tex while Paper 3 is frozen at v0.2.
% Use it as a scratchpad for fully-formal theorem statements, dependency auditing, and proof development.
%
% Canonical theorem-workflow protocol (self-governing):
%   - To add ANY new mathematical claim (definition/remark/proposition/lemma/theorem),
%     COPY the canonical template block below and fill it in.
%   - Do not create ad-hoc formats; the template is the interface.
%
% ============================================================================
% CANONICAL CANDIDATE-RESULT TEMPLATE (COPY FROM HERE FOR EVERY NEW CLAIM)
% ============================================================================
%\begin{comment}
%%%%%%%%%%%%%%%%%%%%%%%%%%%%%%%%%%%%%%%%%%%%%%%%%%%%%%%%%%%%%%
% Candidate Result
%%%%%%%%%%%%%%%%%%%%%%%%%%%%%%%%%%%%%%%%%%%%%%%%%%%%%%%%%%%%%%
% Status:
%   Candidate
%   Definition
%   Remark
%   Proposition
%   Lemma
%   Theorem
%
% ------------------------------------------------------------
% Existence Gate
% ------------------------------------------------------------
% Q1. Exact statement:
% Q2. Minimal definitions required:
% Q3. Counterexample attempts:
% Q4. Classification (circle one): Definition / Remark / Proposition / Lemma / Theorem
%
% Decision:
%
%\begin{theorem}[Working title] % or definition/remark/proposition/lemma
%\label{...}
%\deps{...}
%...
%\end{theorem}
%
% ------------------------------------------------------------
% Dependency Audit
% ------------------------------------------------------------
% Logical dependencies:
% Interface dependencies:
% Downstream users:
%
% ------------------------------------------------------------
% Proof Development
% ------------------------------------------------------------
% Sketch:
% Missing obligations:
% Failed proof attempts:
%
% ------------------------------------------------------------
% Manuscript Promotion Gate
% ------------------------------------------------------------
% Necessity:
% Independence:
% Narrative role:
% Final decision: Stay internal / Promote to main.tex
%%%%%%%%%%%%%%%%%%%%%%%%%%%%%%%%%%%%%%%%%%%%%%%%%%%%%%%%%%%%%%
%\end{comment}
% ============================================================================
% END CANONICAL TEMPLATE
% ============================================================================

\section*{Framework Theorems (v0.3 working draft)}

% ------------------------------------------------------------
% Proposition 1: Typed well-formedness (downgraded from planned Theorem 1)
% ------------------------------------------------------------
\begin{proposition}[Typed well-formedness of structural analysis operators]
\label{prop:typed-well-formedness}
\deps{Definitions~\ref{def:observation-spaces} and~\ref{def:structural-analysis-operator}.}
Let $\mathsf A$ be a structural analysis operator. Then there exists an observation space $Y\in\mathsf{Obs}$ such that $\mathsf A: \Sigma\to Y$ is a deterministic total map on canonical structural objects. In particular, for every canonical structural object $\Sigma$, there exists a unique observable value $y\in Y$ with $y=\mathsf A(\Sigma)$. If $\mathsf A$ is evidence-producing, then for every $\Sigma$ there exists a unique pair $(y,e)$ such that $\mathsf A(\Sigma)=(y,e)$ with $y\in Y$ and $e\in\mathsf{Ev}$.

% Classification (Q4): definitional consequence / typing sanity check.
% Rationale: Definition~\ref{def:structural-analysis-operator} already fixes $\mathsf A$ as a deterministic total function $\Sigma\to Y$ for some $Y\in\mathsf{Obs}$.
\end{proposition}

% Pre-proof checkpoint (answer these before attempting any proof sketch):
%
%  (Q1) What exactly is the statement?
%   - Drafted above (typed existence/uniqueness; evidence-producing case included).
%
%  (Q2) Which definitions are necessary and sufficient?
%   - Necessary (for the statement as written):
%       * Def.~\ref{def:observation-spaces} (to quantify $Y\in\mathsf{Obs}$)
%       * Def.~\ref{def:structural-analysis-operator} (to interpret “$\mathsf A$ is a structural analysis operator”)
%       * (If keeping the evidence clause) Def.~\ref{def:witness-object} only for the *type* $\mathsf{Ev}$.
%       * Canonical structural object $\Sigma$ (from Paper~2) as the domain.
%   - Not currently used by this statement (so should not be dependencies unless we strengthen the claim):
%       * Def.~\ref{def:structural-observable}
%       * Def.~\ref{def:representation-invariant-observable}
%       * Def.~\ref{def:perturbation-family}
%   - Sufficiency check: if Def.~\ref{def:structural-analysis-operator} already includes determinism and a total map $\Sigma\to Y$, then the theorem reduces to typed well-formedness.
%
%  (Q3) Could the statement fail for any taxonomy family? (counterexample attempt)
%   - Evidence-producing operators: OK (still functional), unless evidence is intended to be non-unique.
%   - Partial operators: would break existence for some $\Sigma$ (but are excluded if “operator” means total function).
%   - Nondeterministic analyses: would break uniqueness (excluded by “deterministic” in Def.~\ref{def:structural-analysis-operator}).
%   - Probabilistic outputs: would require $Y$ to be, e.g., a distribution space; allowed only if such $Y\in\mathsf{Obs}$ and the operator is still deterministic as a map into that $Y$.
%   - Multi-output analyses: allowed only if product/tuple spaces are in $\mathsf{Obs}$, or the operator is reframed to return a structured $y$.
%   - Transformational analyses returning a new structural object: allowed only if $\Sigma\in\mathsf{Obs}$ (or a wrapper type is), otherwise this is outside the current interface.
%   - Key fork: if Def.~\ref{def:structural-analysis-operator} is intended to permit partiality or nondeterminism, then Q1 must be tightened (or the definition must be revised).
%
%  (Q4) Is this a theorem, or a consequence of the definitions?
%   - Pending Q2/Q3 resolution. Likely a proposition/remark unless “well-definedness” is strengthened beyond typing.
%
% Proof idea (comments only; only after Q1--Q4 pass):
% - If kept as a result, the “proof” should be a short typing argument.
% - If we want a nontrivial theorem, strengthen the conclusion (e.g., well-definedness relative to $\equiv$, or closure/totality conditions that are not definitional).

% ------------------------------------------------------------
% Theorem 2: Representation invariance
% ------------------------------------------------------------
\begin{theorem}[Representation invariance]
\label{thm:representation-invariance}
\deps{Definition~\ref{def:representation-invariant-observable}; the chosen representation-artifact equivalence $\equiv$ from Paper~2; Definition~\ref{def:structural-analysis-operator}.}
Fix an observation space $Y\in\mathsf{Obs}$ and a representation-invariant structural observable $\mathcal O: \Sigma\to Y$ (with respect to $\equiv$). Define a structural analysis operator $\mathsf A$ by $\mathsf A(\Sigma)\coloneqq \mathcal O(\Sigma)$. Then for all canonical structural objects $\Sigma_1,\Sigma_2$,
\[
\Sigma_1\equiv \Sigma_2 \implies \mathsf A(\Sigma_1)=\mathsf A(\Sigma_2).
\]
If $\mathsf A$ is evidence-producing, the invariance requirement applies to the observable value $y$; evidence objects $e$ need not be invariant across representations.
\end{theorem}

% Existence gate (answer these before attempting any proof sketch):
%
%  (Q1) What exactly is the statement?
%   - Drafted above as: “representation-invariant observables induce representation-invariant analysis results.”
%
%  (Q2) Which definitions are necessary and sufficient?
%   - Necessary:
%       * Def.~\ref{def:representation-invariant-observable} (the hypothesis)
%       * Def.~\ref{def:structural-analysis-operator} (to package $\mathcal O$ as an operator $\mathsf A$)
%       * Equivalence relation $\equiv$ from Paper~2
%   - Not used (unless we strengthen the theorem): witness/evidence definition; perturbation families.
%
%  (Q3) Counterexample attempt:
%   - If $\mathsf A$ is not definitionally tied to a representation-invariant observable (i.e., we allow arbitrary operators), the claim can fail.
%   - If we require invariance of evidence $e$ (not just $y$), the claim can fail for standard witnesses (paths, traces) that depend on naming/order.
%
%  (Q4) Theorem vs. definitional consequence?
%   - Risk: as stated, this is essentially a one-line unfolding of Def.~\ref{def:representation-invariant-observable}.
%   - If we want a substantive theorem, strengthen it to cover a broader class of operators (not just $\mathsf A\coloneqq\mathcal O$) or connect to Paper~2 canonicalization guarantees.
%
% Proof idea (comments only; only after Q1--Q4 pass):
% - Unfold the definition of representation invariance and substitute $\mathsf A=\mathcal O$.
% - Keep evidence out of the equality claim; treat evidence via “checkability” only.

% ------------------------------------------------------------
% Candidate Theorem: Composition
% ------------------------------------------------------------
\begin{comment}
%%%%%%%%%%%%%%%%%%%%%%%%%%%%%%%%%%%%%%%%%%%%%%%%%%%%%%%%%%%%%%
% Candidate Result: Composition
%%%%%%%%%%%%%%%%%%%%%%%%%%%%%%%%%%%%%%%%%%%%%%%%%%%%%%%%%%%%%%
% Status:
%   Candidate
%
% ------------------------------------------------------------
% Existence Gate
% ------------------------------------------------------------
% Q1. Exact statement attempt:
%   Let $Y_1,Y_2\in\mathsf{Obs}$. Let
%   \[
%   \mathsf A_1:\Sigma\to Y_1
%   \qquad\text{and}\qquad
%   \mathsf A_2:Y_1\to Y_2
%   \]
%   be deterministic total maps. If $Y_1$ is an admissible input domain for
%   $\mathsf A_2$ and the composite output space $Y_2$ is admitted as an
%   observation space, then
%   \[
%   \mathsf A_2\circ\mathsf A_1:\Sigma\to Y_2
%   \]
%   is a deterministic total map on canonical structural objects.
%
%   Candidate promotion form, if second-stage maps are admitted as structural
%   analysis operators or as post-observation maps:
%   If $\mathsf A_1$ is an admissible structural analysis into $Y_1$ and
%   $\mathsf A_2$ is an admissible deterministic total post-observation map
%   from $Y_1$ to $Y_2$, then $\mathsf A_2\circ\mathsf A_1$ is an admissible
%   structural analysis into $Y_2$.
%
% Q2. Minimal dependency list:
%   - Canonical structural objects $\Sigma$ as the first-stage domain.
%   - Definition~\ref{def:observation-spaces}, only to require
%     $Y_1,Y_2\in\mathsf{Obs}$.
%   - Definition~\ref{def:structural-analysis-operator}, only for the
%     determinism/totality/admissibility expectations on maps out of $\Sigma$.
%   - A still-unresolved admissibility clause for maps whose domain is an
%     observation space rather than $\Sigma$:
%       * either structural analysis operators are generalized to permit
%         $Y_1\to Y_2$ stages;
%       * or a separate class of deterministic total post-observation maps is
%         admitted;
%       * or the theorem is weakened to a plain typed-composition fact rather
%         than a structural-analysis-operator closure result.
%
% Q3. Counterexample attempts across taxonomy families:
%   - Plain typed maps: if $\mathsf A_1$ and $\mathsf A_2$ are deterministic
%     total maps with matching middle type $Y_1$, composition is typed and
%     total. No additional closure beyond $Y_2\in\mathsf{Obs}$ appears needed.
%   - Structural-analysis-operator taxonomy: if operators are defined only as
%     maps $\Sigma\to Y$, then $\mathsf A_2:Y_1\to Y_2$ is not itself a
%     structural analysis operator unless $Y_1=\Sigma$ or the interface is
%     explicitly extended. The closure risk is therefore admissibility of the
%     second-stage domain, not composition of functions.
%   - Partial analyses: if either stage is partial, the composite can be
%     undefined even when the types match. Totality is required unless partial
%     operators are made explicit.
%   - Nondeterministic analyses: if either stage is nondeterministic, the
%     composite is not a deterministic map without a separate determinization
%     convention. Determinism is required for the current interface.
%   - Evidence-producing analyses: if $\mathsf A_1$ returns $(y_1,e_1)$ while
%     $\mathsf A_2$ expects $y_1$, the middle type does not match. Either the
%     observable projection must be specified, or $\mathsf A_2$ must accept the
%     paired output. No evidence propagation rule is required for the minimal
%     observable-only statement.
%   - Representation-sensitive witnesses: even if observable values compose,
%     evidence objects may fail representation invariance. Therefore any
%     evidence clause should be excluded from the minimal statement unless an
%     admissible evidence interface is separately declared.
%   - Multi-output analyses: if $\mathsf A_1$ returns structured data and
%     $\mathsf A_2$ expects only one component, composition requires an explicit
%     projection or a precisely matching $Y_1$. Do not assume product closure
%     unless the theorem is strengthened to construct paired outputs.
%   - Transformational analyses: if an analysis returns a new structural object
%     and $\mathsf A_2$ is an ordinary structural analysis, composition works
%     only if that returned object type is exactly the canonical structural
%     domain or is explicitly admitted as an observation/input domain. This is
%     another domain-admissibility issue.
%   - Probabilistic-valued analyses: composition is acceptable only when the
%     probabilistic object is treated as an element of $Y_1$ and
%     $\mathsf A_2$ is a deterministic total map on that object. Sampling-based
%     second stages would violate determinism.
%
% Q4. Classification risk: theorem vs proposition:
%   - As a plain statement about deterministic total typed maps, this is at
%     most a proposition or remark: ordinary function composition supplies the
%     result.
%   - It becomes theorem-worthy only if the framework first defines a
%     nontrivial admissibility/closure rule saying when post-observation maps
%     $Y_1\to Y_2$ preserve structural-analysis admissibility.
%   - Current risk: calling it a theorem would overstate a definitional or
%     typing consequence unless the missing admissibility interface is added.
%
% Decision:
%   Keep as Candidate. Do not prove yet. Next action is to determine whether
%   Paper 3 wants (i) a post-observation-map admissibility definition, (ii) a
%   weakened typed-composition proposition, or (iii) no promoted result.
%
% ------------------------------------------------------------
% Dependency Audit
% ------------------------------------------------------------
% Logical dependencies:
%   Canonical structural objects; deterministic total functions; admitted
%   observation spaces; unresolved post-observation admissibility.
% Interface dependencies:
%   Matching middle type $Y_1$; $Y_2\in\mathsf{Obs}$; either admissibility of
%   $Y_1$ as a second-stage input or explicit weakening to plain composition.
% Downstream users:
%   Any later result that chains analyses without re-reading the original
%   structural object.
%
% ------------------------------------------------------------
% Proof Development
% ------------------------------------------------------------
% Sketch:
%   No proof yet.
% Missing obligations:
%   Determine whether the framework admits maps out of observation spaces and
%   whether evidence-producing outputs are excluded, projected, or typed as the
%   actual middle object.
% Failed proof attempts:
%   None yet; proof intentionally deferred.
%
% ------------------------------------------------------------
% Manuscript Promotion Gate
% ------------------------------------------------------------
% Necessity:
%   Pending downstream dependency audit.
% Independence:
%   Pending decision on whether this is more than typed function composition.
% Narrative role:
%   Could justify analysis pipelines only after the admissibility interface is
%   fixed.
% Final decision: Stay internal.
%%%%%%%%%%%%%%%%%%%%%%%%%%%%%%%%%%%%%%%%%%%%%%%%%%%%%%%%%%%%%%
\end{comment}

% ------------------------------------------------------------
% Theorem 4: Perturbation compatibility
% ------------------------------------------------------------
\begin{theorem}[Perturbation compatibility]
\label{thm:perturbation-compatibility}
\deps{Definition~\ref{def:perturbation-family}; the admissibility conditions for analyses and perturbations.}
% TODO (fully formal): define the compatibility notion (commutation, monotonicity, Lipschitz/stability, simulation, etc.).
% TODO: quantify over $\delta\in\Delta$ and relate $\mathsf A(\Sigma)$ to $\mathsf A(\delta(\Sigma))$.
\end{theorem}

% Theorem viability pass (answer before attempting any proof sketch):
%  (1) Meaningful?  [TODO]
%  (2) Nontrivial?  [TODO: identify a hypothesis whose removal breaks the claim]
%  (3) Needed?     [TODO: list downstream theorems/lemmas that require it]
%  (4) Likely true under current definitions? [TODO]
%
% Proof idea (comments only):
% - Full commutativity is often too strong; consider a weaker relation (e.g., a refinement preorder on outputs, a bound, or an explicit “delta observable”).
% - If the output space is quantitative, compatibility may become a stability inequality.

% End of file.
